% CREATED BY DAVID FRISK, 2016
\oneLineTitle\\
% \oneLineSubtitle\\
NIKOLAOS ALEXANDRIS\\
Department of Computer Science and Engineering\\
Chalmers University of Technology and University of Gothenburg

\thispagestyle{plain}			% Supress header
\section*{Abstract}

Compiler construction is a research area that is constantly being developed, as
the need for faster and more correct programming languages and implementations
is steadily growing. While much work has been done to ensure compilers output
efficient and semantically intact code, the complexity and size of the current
optimising compilers lead to bugs that can be hard to find and diagnose.
Situations where correctness is critical has therefore led to the development of
formally verified compilers, that guarantee the correctness of the generated
code. The price paid for developing verified compilers is that every step in the
compilation process must be carefully implemented and formalised in order for
the verification system to accept it, which makes optimisation pipelines much
harder to implement.

This project's work is based on the PureCake verified compiler for pure, lazy
functional languages. PureCake is an extension of the CakeML project, which
itself is a verified compiler for strict functional languages, both developed
using the HOL4 formal verification system. The PureCake compiler is already a
complete verified compiler, but it lacks the optimisations that would make its
generated code as efficient as the generated code of optimising non-verified
compilers for lazy languages such as GHC. This project's goal is to develop and
verify critical missing optimisations, both in PureCake's optimisation pipeline
and in the CakeML's runtime system, in order to close the gap between verified
and non-verified compilers for lazy languages in terms of runtime efficiency.

In this thesis, we take on the task of improving the performance of PureCake's
generated code, by optimising the handling of thunks by the PureCake and CakeML
compilers. We first target the representation of thunks inside the compilation
pipeline, by implementing optimisations that try to eliminate the generation of
as many thunks as possible. Then, we target the runtime representation of the
remaining thunks so that they have a minimal effect on the performance of the
executable. This work will help bring PureCake up to par with unverified
compilers in terms of the generated executable's performance, while also making
sure that the soundness of the compilation process remains intact.

% KEYWORDS (MAXIMUM 10 WORDS)
\vfill
Keywords: Computer, science, computer science, engineering, project, thesis.

\newpage				% Create empty back of side
\thispagestyle{empty}
\mbox{}
